\subsection{SYLOW SUBGROUPS}

DEFINITION 10. Let G be a finite group. A Sylow p-subgroup of G is any subgroup P of G satisfying the two following conditions:

\begin{enumerate}
\def\labelenumi{(\alph{enumi})}
\item
  P is a p-group.
\item
  (G:P) is not a multiple of p.
\end{enumerate}

If the order of G is written in the form \(p^{r}m\) , where m is not a multiple of p, conditions (a) and (b) are equivalent to \(\operatorname{Card}(P) = p^{r}\) .

Examples. (1) In the group \(\mathfrak{S}_p\) let \(\zeta\) be a cycle of order \(p\) . The subgroup generated by \(\zeta\) is a Sylow \(p\) -subgroup of \(\mathfrak{S}_p\) for \(p\) does not divide \((p - 1)!\) .

\begin{enumerate}
\def\labelenumi{(\arabic{enumi})}
\setcounter{enumi}{1}
\tightlist
\item
  *Let \(k\) be a finite field of characteristic \(p\) and let \(n\) be a positive integer. The strict triangular group \(T_1(n, k)\) is a Sylow \(p\) -subgroup of the group \(\mathbf{GL}(n, k)\) .*
\end{enumerate}

THEOREM 2. Every finite group contains a Sylow p-subgroup.

The proof depends on the following lemma.

Lemma. Let \(n = p^r m\) , where \(m\) is an integer which is not a multiple of \(p\) . Then

\[
\binom {n} {p ^ {r}} \not \equiv 0 \pmod {p}.
\]

Let S be a group of order \(p^r\) (for example \(\mathbf{Z} / p^r\mathbf{Z}\) ) and T a set with \(m\) elements. Write \(\mathbf{X} = \mathbf{S} \times \mathbf{T}\) and let E be the set of subsets of X with \(p^r\) elements. Then \(\operatorname{Card}(\mathbf{X}) = n\) , whence \(\operatorname{Card}(\mathbf{E}) = \binom{n}{p^r}\) (Set Theory, III, § 5, no. 8, Corollary 1 to Proposition 11). Let S operate on X by \(s.(x,y) = (sx,y)(s,x \in \mathbf{S},y \in \mathbf{T})\) and consider the canonical extension of this operation to E. In the notation of no. 5, Proposition 11, the set \(\mathbf{E}^{\mathbf{S}}\) is the set of orbits of X, that is the set of subsets \(Y \subset X\) of the form \(S \times \{t\}, t \in T\) , whence \(\operatorname{Card}(\mathbf{E}^{\mathbf{S}}) = m\) . By no. 5, Proposition 11,

\[
\binom {n} {p ^ {r}} = \operatorname{Card} (\mathrm{E}) \equiv \operatorname{Card} (\mathrm{E} ^ {\mathrm{S}}) = m \not \equiv 0 \pmod {p},
\]

which proves the lemma.

We now prove the theorem. Let G be a finite group and n its order; we write \(n = p^{r}m\) , where m is not a multiple of p.~Let E be the set of subsets of G with \(p^{r}\) elements. Then

\[
\operatorname{Card} (\mathbf {E}) = \binom{n}{p ^ {r}};
\]

whence, by virtue of the lemma, \(\operatorname{Card}(\mathbf{E}) \not\equiv 0 (\operatorname{mod.} p)\) . Consider the extension to E of the operation of G on itself by left translation. There exists \(X \in E\) whose orbit has non-zero cardinal mod. p.~If \(H_{X}\) denotes the stabilizer of X, then \((G:H_{X}) \not\equiv 0 (\text{mod. } p)\) , which means that \(p^{r}\) divides \(\text{Card}(H_{X})\) . But \(H_{X}\) consists of the \(s \in G\) such that \(sX = X\) ; if \(x \in X\) , then \(H_{X} \subset X.x^{-1}\) , whence \(\text{Card}(H_{X}) \leqslant \text{Card}(X) = p^{r}\) . Hence \(\text{Card}(H_{X}) = p^{r}\) .

COROLLARY. If the order of G is divisible by p, the group G contains an element of order p.

By virtue of Theorem 2, this is reduced to the case where G is a p-group \(\neq\{e\}\) ; if \(x\in G\) is different from e, the cyclic group generated by x is then of order \(p^{n}\) with \(n\geqslant1\) and it therefore contains a subgroup of order p.

Remark. For every prime number q dividing \(\operatorname{Card}(\mathbf{G})\) , let \(P_{q}\) be a Sylow q-subgroup of G. Then the subgroup H of G generated by the \(P_{q}\) is of order a multiple of \(\operatorname{Card}(\mathbf{P}_{q})\) for each q and of order a divisor of \(\operatorname{Card}(\mathbf{G})\) , hence it is equal to G.

THEOREM 3. Let G be a finite group.

\begin{enumerate}
\def\labelenumi{(\alph{enumi})}
\item
  The Sylow \(p\) -subgroups of \(\mathbf{G}\) are conjugate to one another. Their number is congruent to 1 mod. \(p\) .
\item
  Every subgroup of G which is a p-group is contained in a Sylow p-subgroup.
\end{enumerate}

Let P be a Sylow p-subgroup of G (Theorem 2) and let H be a p-subgroup of G. Let E = G/P and consider the operation of H on G/P. As \(\text{Card}(E) \not\equiv 0 \mod p\) , Proposition 11 of no. 5, shows that there exists \(x \in G/P\) such that hx = x for all \(h \in H\) . If g is a representative of x in G, this means that \(H \subset gPg^{-1}\) , whence assertion (b).

If H is a Sylow p-subgroup, then \(\operatorname{Card}(\mathrm{H}) = \operatorname{Card}(\mathrm{P}) = \operatorname{Card}(g\mathrm{Pg}^{-1})\) , whence \(H = gPg^{-1}\) , which proves the first assertion of (a).

We now prove the second assertion of (a). Let \(\mathcal{S}\) be the set of Sylow \(p\) -subgroups of G and let P operate on \(\mathcal{S}\) by inner automorphisms. The element \(\mathrm{P} \in \mathcal{S}\) is a fixed point under this operation, we show that it is the only one. Let \(\mathrm{Q} \in \mathcal{S}\) be a fixed point; Q is a Sylow subgroup of G normalized by P and hence P is contained in the normalizer N of Q. The groups P and Q are Sylow \(p\) -subgroups of N; hence there exists \(n \in N\) such that \(\mathrm{P} = n\mathrm{Q}n^{-1} = \mathrm{Q}\) . By no. 5, Proposition 11, \(\operatorname{Card}(\mathcal{S}) \equiv \operatorname{Card}(\mathcal{S}^{\mathrm{P}}) = 1 (\text{mod. } p)\) .

COROLLARY 1. Let P be a Sylow p-subgroup of G, let N be its normalizer in G and let M be a subgroup of G containing N. The normalizer of M in G is equal to M.

Let \(s \in G\) be such that \(sMs^{-1} = M\) . The subgroup \(sPs^{-1}\) of M is a Sylow \(p\) -subgroup of M. There thus exists \(t \in M\) such that \(sPs^{-1} = tPt^{-1}\) ; then \(t^{-1}s \in N\) , whence \(s \in tN \subset M\) .

COROLLARY 2. Let \(f: \mathbf{G}_{1} \to \mathbf{G}_{2}\) be a homomorphism of finite groups. For every Sylow \(p\) -subgroup \(\mathbf{P}_{1}\) of \(\mathbf{G}_{1}\) there exists a Sylow \(p\) -subgroup \(\mathbf{P}_{2}\) of \(\mathbf{G}_{2}\) such that \(f(\mathbf{P}_{1}) \subset \mathbf{P}_{2}\) .

This follows from Theorem 3 (b) applied to the subgroup \(f(\mathbf{P}_1)\) of \(\mathbf{G}_2\) .

COROLLARY 3. (a) Let H be a subgroup of G. For every Sylow p-subgroup P of H there exists a Sylow p-subgroup Q of G such that P = Q ∩ H.

\begin{enumerate}
\def\labelenumi{(\alph{enumi})}
\setcounter{enumi}{1}
\item
  Conversely, if \(\mathbf{Q}\) is a Sylow \(p\) -subgroup of \(\mathbf{G}\) and \(\mathbf{H}\) is normal in \(\mathbf{G}\) , the group \(\mathbf{Q} \cap \mathbf{H}\) is a Sylow \(p\) -subgroup of \(\mathbf{H}\) .
\item
  (a) The \(p\) -group \(\mathbf{P}\) is contained in a Sylow \(p\) -subgroup \(\mathbf{Q}\) of \(\mathbf{G}\) and \(\mathbf{Q} \cap \mathbf{H}\) is a \(p\) -subgroup of \(\mathbf{H}\) containing \(\mathbf{P}\) and is hence equal to \(\mathbf{P}\) .
\item
  (b) Let \(\mathbf{P}'\) be a Sylow \(p\) -subgroup of H. There exists an element \(g \in G\) such that \(g\mathbf{P}'g^{-1} \subset \mathbf{Q}\) . As H is normal, \(\mathbf{P} = g\mathbf{P}'g^{-1}\) is contained in H and hence in \(\mathbf{Q} \cap \mathbf{H}\) . As \(\mathbf{Q} \cap \mathbf{H}\) is a \(p\) -subgroup of H and P is a Sylow \(p\) -subgroup of H, \(\mathbf{P} = \mathbf{Q} \cap \mathbf{H}\) .
\end{enumerate}

COROLLARY 4. Let N be a normal subgroup of G. The image in G/N of a Sylow p-subgroup of G is a Sylow p-subgroup of G/N and every Sylow p-subgroup of G/N is obtained in this way.

Let \(G' = G/N\) and \(P'\) be the image in \(G'\) of a Sylow \(p\) -subgroup \(P\) of \(G\) . The group \(G\) operates transitively on \(G'/P'\) and hence \(G'/P'\) is equipotent to \(G/S\) , where \(S\) is a subgroup of \(G\) containing \(P\) . Therefore \((G':P')\) divides \((G:P)\) , is thus not a multiple of \(p\) and the \(p\) -group \(P'\) is a Sylow \(p\) -subgroup of \(G'\) . Let \(Q'\) be another Sylow \(p\) -subgroup of \(G'\) ; then \(Q' = g'P'g'^{-1}\) for some \(g' \in G'\) ; if \(g \in G\) is a representative of \(g'\) , the group \(Q'\) is the image of \(Q = gPg^{-1}\) .

